\subsection{Transforming a Graph: Order of Operations} 
Let's summarize what we've found.
\begin{itemize}
\item Addition and subtraction correspond to \makeblank[1in].
\item Multiplication comes in two parts:
\begin{itemize}
\item Multiplication by a \makeblank[.75in] number corresponds to \makeblank[1in].
\item If the number is \makeblank[.75in], there is also a \makeblank[1in].
\end{itemize}
\item Anything \emph{inside} the function has a \makeblank[1.5in] effect.
\item Anything \emph{outside} the function has a \makeblank[1.5in] effect.
\item \makeblank[1.5in] effects are \makeblank[1.5in].
\end{itemize}
What about order?
\begin{itemize}
\item \emph{Outside} the function, we follow order of operation: \makeblank[1in] before \makeblank[1in].
\item \emph{Inside} the function, we \emph{reverse} order of operations (remember: everything inside the function is \makeblank[1in]!), so \makeblank[1in] before \makeblank[1in].
\end{itemize}
When we work, we will usually work on the \makeblank rather than on the picture directly.
\begin{example}
Below is the graph of $y=f(x)$. Sketch the graph of $y=-2f(\frac{2}{3}x+5)+3$
\begin{icpic}[x=.5\linespace,y=.5\linespace]
\useasboundingbox (-13,-8) rectangle (13,8);
\setgraph{13}{13}{8}{8}
\GraphSmall
\draw[graphend] (-3,2) -- (1,5) -- (7,-1);
\draw[graphend] (1,5) node[anchor=west]{$y=f(x)$};
%\draw[thick,red] (-12,-1) -- (-6,-7) -- (3,1);
\end{icpic}
\steps[0]
\begin{tabular}{r<{\numstep}@{\;}l>{$}l<{$}>{\blankpair}c>{\blankpair}c>{\blankpair}c}
&Starting Point&{}&&&\\
&Horizontal translation&x\rightarrow\makeblanksmall&&&\\
&Horiz. scaling/reflection&x\rightarrow\makeblanksmall&&&\\
&Vert. scaling/reflection&y\rightarrow\makeblanksmall&&&\\
&Vertical translation&y\rightarrow\makeblanksmall&&&
\end{tabular}
\end{example}